\documentclass[11pt,a4paper]{article}

\usepackage[utf8]{inputenc}
\usepackage[T1]{fontenc}
\usepackage[spanish,es-noshorthands]{babel}
\usepackage[margin=2.2cm]{geometry}
\usepackage{booktabs}
\usepackage{tikz}
\usetikzlibrary{arrows.meta,positioning,fit}
\usepackage{float}

\title{\vspace{-1.5cm}Simulador MIPS con \emph{pipeline} de 5 etapas\\
\large Diagramas del proyecto (Parte 2)}
\author{}
\date{\today}

\begin{document}
\maketitle
\vspace{-1.2cm}

\section{Diagrama del sistema}

\begin{figure}[H]
\centering
\begin{tikzpicture}[
  font=\small,
  etapa/.style={draw,rounded corners,minimum width=4.2cm,minimum height=1.05cm,align=center,fill=violet!8},
  reg/.style={draw,minimum width=4.2cm,minimum height=0.45cm,align=center,fill=black!6,font=\scriptsize},
  rec/.style={draw,rounded corners,minimum width=2.4cm,minimum height=1.05cm,align=center,fill=teal!12,font=\scriptsize},
  test/.style={draw,rounded corners,minimum width=3.1cm,minimum height=0.95cm,align=center,fill=orange!13,font=\scriptsize},
  fl/.style={-{Latex[length=2mm]}},
  bi/.style={{Latex[length=2mm]}-{Latex[length=2mm]}}
]

\node[etapa] (if) {\textbf{1. IF --- Fetch + PC}\\[1pt]{\scriptsize \texttt{stage\_if.c}}};
\node[reg,below=3.5mm of if] (r1) {registro IF/ID};
\node[etapa,below=3.5mm of r1] (id) {\textbf{2. ID --- Decode}\\[1pt]{\scriptsize \texttt{stage\_id.c}}};
\node[reg,below=3.5mm of id] (r2) {registro ID/EX};
\node[etapa,below=3.5mm of r2] (ex) {\textbf{3. EX --- ALU}\\[1pt]{\scriptsize \texttt{stage\_ex.c}}};
\node[reg,below=3.5mm of ex] (r3) {registro EX/MEM};
\node[etapa,below=3.5mm of r3] (mem) {\textbf{4. MEM --- Memory}\\[1pt]{\scriptsize \texttt{stage\_mem.c}}};
\node[reg,below=3.5mm of mem] (r4) {registro MEM/WB};
\node[etapa,below=3.5mm of r4] (wb) {\textbf{5. WB --- Write back}\\[1pt]{\scriptsize \texttt{stage\_wb.c}}};

\foreach \a/\b in {if/r1,r1/id,id/r2,r2/ex,ex/r3,r3/mem,mem/r4,r4/wb}
  \draw[fl] (\a) -- (\b);

\node[rec,left=1.5cm of if]  (imem) {\textbf{Mem. instrucciones}\\[1pt]\texttt{imem.bin}};
\node[rec,left=1.5cm of id]  (rf)   {\textbf{Banco de registros}\\[1pt]\texttt{regs.bin}};
\node[rec,left=1.5cm of mem] (dmem) {\textbf{Mem. de datos}\\[1pt]\texttt{dmem.bin}};

\draw[fl] (imem) -- (if);
\draw[fl] (rf)   -- (id);
\draw[bi] (dmem) -- (mem);
\draw[fl] (wb.west) -- ++(-0.75,0) |- ([yshift=-3.5mm]rf.east);

\node[test,right=1.5cm of if]  (tif)  {\textbf{\texttt{test\_if}}\\[1pt]PC+4, saltos, \texttt{jr}};
\node[test,right=1.5cm of id]  (tid)  {\textbf{\texttt{test\_id}}\\[1pt]control, extensión};
\node[test,right=1.5cm of ex]  (tex)  {\textbf{\texttt{test\_ex}}\\[1pt]6 operaciones, zero};
\node[test,right=1.5cm of mem] (tmem) {\textbf{\texttt{test\_mem}}\\[1pt]\texttt{lw}, \texttt{sw}, rangos};
\node[test,right=1.5cm of wb]  (twb)  {\textbf{\texttt{test\_wb}}\\[1pt]\texttt{\$0}, \texttt{memToReg}};

\foreach \a/\b in {if/tif,id/tid,ex/tex,mem/tmem,wb/twb}
  \draw[bi] (\a) -- (\b);

\node[test,below=7mm of wb,minimum width=8.5cm] (sys)
  {\textbf{\texttt{test\_integration}}\\[1pt]vectores S1--S8: corre el programa y compara el estado final};
\draw[fl] (wb) -- (sys);

\end{tikzpicture}
\caption{Módulos del simulador, recursos de estado cargados desde \texttt{.bin},
módulos de prueba y conexiones direccionadas.}
\end{figure}

\clearpage
\section{Diagrama de subfunciones de cada función}

\begin{figure}[H]
\centering
\begin{tikzpicture}[
  font=\small,
  sub/.style={draw,rounded corners=2pt,fill=violet!8,minimum width=2.2cm,minimum height=0.62cm,
              font=\scriptsize,align=center},
  cont/.style={draw,rounded corners=4pt,fill=black!4}
]

\draw[cont] (0,0.45) rectangle (12.2,-1.35);
\node[anchor=west,font=\bfseries\small] at (0.25,0) {1. IF --- Fetch + PC};
\node[anchor=east,font=\tiny] at (11.95,0)
  {entra: \texttt{pc[31:0]} $\cdot$ sale: \texttt{instr[31:0]}, \texttt{pcPlus4[31:0]}};
\node[sub] at (1.3,-0.85) {\texttt{selectNextPC}};
\node[sub] at (3.7,-0.85) {\texttt{checkAlign}};
\node[sub] at (6.1,-0.85) {\texttt{fetchInstr}};
\node[sub] at (8.5,-0.85) {\texttt{incrementPC}};

\draw[cont] (0,-1.95) rectangle (12.2,-3.75);
\node[anchor=west,font=\bfseries\small] at (0.25,-2.4) {2. ID --- Decode};
\node[anchor=east,font=\tiny] at (11.95,-2.4)
  {entra: \texttt{instr}, \texttt{pcPlus4} $\cdot$ sale: \texttt{rsVal}, \texttt{rtVal}, \texttt{imm32}, control};
\node[sub] at (1.3,-3.25) {\texttt{splitFields}};
\node[sub] at (3.7,-3.25) {\texttt{controlUnit}};
\node[sub] at (6.1,-3.25) {\texttt{readRegs}};
\node[sub] at (8.5,-3.25) {\texttt{signExtend}};
\node[sub] at (10.9,-3.25) {\texttt{calcTargets}};

\draw[cont] (0,-4.35) rectangle (12.2,-6.15);
\node[anchor=west,font=\bfseries\small] at (0.25,-4.8) {3. EX --- ALU};
\node[anchor=east,font=\tiny] at (11.95,-4.8)
  {entra: \texttt{a}, \texttt{b}, \texttt{aluOp[1:0]}, \texttt{funct[5:0]} $\cdot$ sale: \texttt{result[31:0]}, \texttt{zero}};
\node[sub] at (1.3,-5.65) {\texttt{aluControl}};
\node[sub] at (3.7,-5.65) {\texttt{selectOperB}};
\node[sub] at (6.1,-5.65) {\texttt{execute}};
\node[sub] at (8.5,-5.65) {\texttt{zeroFlag}};
\node[sub] at (10.9,-5.65) {\texttt{branchTaken}};

\draw[cont] (0,-6.75) rectangle (12.2,-8.55);
\node[anchor=west,font=\bfseries\small] at (0.25,-7.2) {4. MEM};
\node[anchor=east,font=\tiny] at (11.95,-7.2)
  {entra: \texttt{address}, \texttt{writeData}, \texttt{memRead/Write} $\cdot$ sale: \texttt{readData[31:0]}};
\node[sub] at (1.3,-8.05) {\texttt{mapAddress}};
\node[sub] at (3.7,-8.05) {\texttt{checkBounds}};
\node[sub] at (6.1,-8.05) {\texttt{readWord}};
\node[sub] at (8.5,-8.05) {\texttt{writeWord}};

\draw[cont] (0,-9.15) rectangle (12.2,-10.95);
\node[anchor=west,font=\bfseries\small] at (0.25,-9.6) {5. WB};
\node[anchor=east,font=\tiny] at (11.95,-9.6)
  {entra: \texttt{readData}, \texttt{aluResult}, \texttt{memToReg} $\cdot$ sale: banco de registros};
\node[sub] at (1.3,-10.45) {\texttt{selectData}};
\node[sub] at (3.7,-10.45) {\texttt{selectDest}};
\node[sub] at (6.1,-10.45) {\texttt{writeRegFile}};

\end{tikzpicture}
\caption{Subfunciones de cada una de las cinco funciones, con sus entradas y salidas.
Cada nombre corresponde a una función estática dentro del archivo \texttt{.c} de la etapa.}
\end{figure}

\clearpage
\section{Correspondencia entre el diseño y el código}

\begin{table}[H]
\centering\small
\begin{tabular}{@{}lll@{}}
\toprule
Módulo del diagrama & Archivo & Función principal \\
\midrule
1. Instruction Fetch + PC & \texttt{src/stage\_if.c}  & \texttt{stage\_if()} \\
2. Instruction Decode     & \texttt{src/stage\_id.c}  & \texttt{stage\_id()} \\
3. ALU (Execute)          & \texttt{src/stage\_ex.c}  & \texttt{stage\_ex()} \\
4. Memory                 & \texttt{src/stage\_mem.c} & \texttt{stage\_mem()} \\
5. Write Back             & \texttt{src/stage\_wb.c}  & \texttt{stage\_wb()} \\
\midrule
Registros de \emph{pipeline}, riesgos y \emph{flush} & \texttt{src/pipeline.c} & \texttt{pipeline\_step()} \\
Memorias y banco de registros (\texttt{.bin})        & \texttt{src/memory.c}   & \texttt{mem\_load\_*()} \\
Módulos de prueba por etapa                          & \texttt{tests/test\_*.c} & un binario por etapa \\
Test del sistema (vectores S1--S8)                   & \texttt{tests/test\_integration.c} & --- \\
\bottomrule
\end{tabular}
\end{table}

\end{document}
